\chapter{Introduction}\label{ch:intro}


Quantum computers have crossed from laboratory curiosities to programmable, cloud-accessible machines with more than a hundred qubits. But no quantum processor today can complete a computation of practical value without enormous classical computation around it. Every circuit is designed, compiled, simulated and debugged on classical hardware before it runs. Classical loops optimize its parameters, and classical decoders post-process, mitigate and, eventually, error-correct its noisy outputs. Classical infrastructure schedules, transports and secures the entire workflow (\cref{fig:intro:lifecycle}).

This dissertation starts from one observation: that classical side is not a temporary scaffold to be removed once the hardware matures. It is a permanent, load-bearing part of the architecture whose design determines how much useful work a quantum processor can deliver. The thesis defended here has two parts. First, a \emph{co-designed stack} will deliver useful quantum computation in the coming decade: classical high-performance computing (HPC) surrounds the quantum processing unit (QPU) at every layer. Second, each layer poses concrete research problems whose solutions compound.

% alt: A left-to-right chain of five rounded boxes for one quantum computation (design and compile; simulate and debug; optimize parameters; execute on the QPU; post-process, mitigate and decode), with a return arrow from the last box to the optimizer for the classical loop, and a band under the chain for the classical infrastructure that schedules, transports and secures every step.
\begin{figure}[htbp]
\centering
\begin{tikzpicture}[node distance=0.45cm and 0.3cm,
  stp/.style={draw,rounded corners=4pt,fill=cbBlue!15,text width=2.25cm,align=center,font=\footnotesize\linespread{1}\selectfont,inner ysep=5pt,minimum height=1.1cm},
  qstp/.style={stp,fill=cbOrange!25,thick},
  arr/.style={-{Stealth[length=2mm]},thick}]
  \node[stp] (design) {design and compile};
  \node[stp,right=of design] (sim) {simulate and debug};
  \node[stp,right=of sim] (opt) {optimize parameters};
  \node[qstp,right=of opt] (qpu) {execute on the QPU};
  \node[stp,right=of qpu] (post) {post-process, mitigate and decode};
  \draw[arr] (design) -- (sim);
  \draw[arr] (sim) -- (opt);
  \draw[arr] (opt) -- (qpu);
  \draw[arr] (qpu) -- (post);
  \draw[arr,cbGray] (post.north) -- ++(0,0.5) -| node[pos=0.25,above,font=\footnotesize] {classical optimization loop} (opt.north);
  \coordinate (bl) at ([yshift=-1.35cm]design.south west);
  \coordinate (tr) at ([yshift=-0.6cm]post.south east);
  \draw[rounded corners=4pt,fill=cbGray!12] (bl) rectangle (tr);
  \node[font=\footnotesize,align=center] at ($(bl)!0.5!(tr)$) {classical infrastructure: scheduling, transport and security of every step};
  \draw[arr,cbGray] (design.south |- tr) -- (design.south);
  \draw[arr,cbGray] (sim.south |- tr) -- (sim.south);
  \draw[arr,cbGray] (opt.south |- tr) -- (opt.south);
  \draw[arr,cbGray] (qpu.south |- tr) -- (qpu.south);
  \draw[arr,cbGray] (post.south |- tr) -- (post.south);
\end{tikzpicture}
\caption{Classical computation around one quantum computation. Every step except execution runs on classical hardware (blue): classical loops optimize the parameters, and classical decoders post-process, mitigate and error-correct the outputs. The band marks the classical infrastructure that schedules, transports and secures every step. Schematic.}
\label{fig:intro:lifecycle}
\end{figure}

\section{Motivation}\label{sec:intro:motivation}

Four gaps motivate the work. The first is a \emph{scale} gap in simulation. Exact state-vector simulation of an $n$-qubit circuit requires $2^{n}$ complex amplitudes, and every additional qubit doubles the memory and time. A 32-qubit state vector in single precision occupies \SI{32}{\gibi\byte}, which fits on one data-center GPU. A 42-qubit state vector requires tens of terabytes, which only a large fraction of a leadership-class supercomputer can hold.

Researchers nevertheless need simulation at exactly this scale. First, it is the regime where today's hardware produces results that can neither be trusted nor fully verified. Second, the algorithms of \cref{ch:nisq,ch:qml} were developed and validated in simulation before they ran on hardware.

GPU simulators are mature. The practical obstacle is the friction of using them. Circuits written in the dominant Python framework, \qiskit, cannot be handed to GPU-native simulators such as \cudaq{} without a rewrite. Deployment across hundreds of nodes of a Slurm-managed system, with containers and MPI, is a specialist task. \Cref{ch:qgear} addresses this gap.

The second is a \emph{data} gap at the classical--quantum interface. Most proposed quantum machine-learning models spend more resources on loading classical data into a quantum state and reading results back than on computing. Angle encodings use one qubit per feature. Amplitude encodings require circuits of depth exponential in the number of qubits. Readout of a single expectation value to precision $\epsilon$ requires $\bigO{1/\epsilon^{2}}$ shots.

Encodings such as \qcrank{} change this arithmetic. They load many data values in parallel into a small register and read out many expectation values in parallel from the same circuit. They make shot-efficient, hardware-compatible learning primitives possible. \Cref{ch:qml} builds four such primitives.

The third is a \emph{noise} gap in algorithm design. Variational algorithms such as the quantum approximate optimization algorithm (QAOA) treat the quantum processor as a black box. A classical optimizer tunes its parameters from noisy expectation-value estimates. Current superconducting devices have two-qubit gate errors near $10^{-3}$ and coherence times of a few hundred microseconds. There, the optimization loop spends most of its shot budget to fight noise. It converges slowly and often fails to find the ground state that a classical solver would find in milliseconds.

The problem instances, however, are not black boxes. A quadratic unconstrained binary optimization (QUBO) matrix carries structure that can be mapped directly to circuit parameters. The device carries calibration data that can guide the placement of logical qubits on physical qubits. \Cref{ch:nisq} develops algorithms that exploit both kinds of structure.

The fourth is a \emph{trust} gap. Users access quantum processors through shared clouds. Error-correcting codes are public. Future quantum computers will attack the classical infrastructure that transports jobs, calibration data and results. Fault-tolerant quantum computing is usually analysed against a benign noise model. A co-located adversary who can inject faults into a shared device is a qualitatively different threat, and a fixed public encoder is a fixed target.

The cryptography that protects the classical side of the stack must also migrate to the post-quantum standards that NIST finalized in 2024. \Cref{ch:qec,ch:pqc} address both sides of this gap. Between them, \cref{ch:synthesis} asks what fault-tolerant circuits and end-to-end quantum simulations will actually cost.

\section{Thesis Statement and Research Questions}\label{sec:intro:thesis}

\begin{quote}
Useful quantum computation in the coming decade will be delivered not by a quantum processor alone but by a co-designed stack in which classical high-performance computing surrounds the QPU at every layer:
\begin{itemize}[nosep,leftmargin=1.5em]
\item GPU-accelerated simulation to develop, verify and calibrate circuits at scales the hardware cannot yet reach;
\item structured, vectorized data encodings that make the classical--quantum data interface efficient and learnable;
\item hardware-aware algorithm design that replaces black-box variational optimization with structure derived from the problem and the device;
\item AI- and HPC-driven synthesis that compiles toward fault-tolerant resource budgets; and
\item security-aware error correction and post-quantum cryptography so that the same stack can be trusted when it runs on shared cloud hardware in the post-quantum era.
\end{itemize}
\end{quote}

Five research questions, one per layer of the stack, organize the dissertation. They are answered in \cref{ch:qgear,ch:qml,ch:nisq,ch:synthesis,ch:qec,ch:pqc} and revisited in \cref{ch:conclusion}.

\begin{description}[style=nextline,leftmargin=1.5em]
\item[RQ1 (Scale).] How far can classical HPC push exact circuit simulation, and how do we make that capability usable without rewriting quantum programs?
\item[RQ2 (Data).] How should classical data enter and leave a quantum processor so that learning models are shot-efficient, hardware-compatible and trainable?
\item[RQ3 (Noise).] How can optimization algorithms on noisy intermediate-scale quantum (NISQ) devices exploit problem and hardware structure instead of black-box parameter search?
\item[RQ4 (Synthesis).] Can learned, HPC-scale synthesis produce circuits that respect fault-tolerant resource budgets, and what does honest cost accounting say about end-to-end quantum advantage?
\item[RQ5 (Trust).] How do we make error-corrected computation on shared cloud QPUs auditable and secure, and how does that connect to post-quantum cryptography standards?
\end{description}

\Cref{fig:intro:stack} shows the five layers, the chapter for each, and the flow of artifacts between them. Simulation results and containers from the bottom layer feed every layer above it. Hardware-aware compilation techniques from the encoding and algorithm layers feed the synthesis layer. The error-correction audit consumes the randomness and key-management services of the security layer. The QPU is deliberately in the middle, not on top of a hierarchy, with the classical layers around it.

% alt: A block diagram with a quantum processor in the centre surrounded by five concentric classical layers, each labelled with a research question and the chapter that addresses it, and arrows showing that simulation artifacts flow upward while randomness and keys flow from the security layer to error correction.
\begin{figure}[htbp]
\centering
\begin{tikzpicture}[node distance=0.35cm,
  layer/.style={draw,rounded corners=4pt,text width=10.2cm,align=center,font=\small\linespread{1}\selectfont,inner ysep=6pt},
  arr/.style={-{Stealth[length=2mm]},thick}]
  \node[layer,fill=cbGray!12] (L5) {\textbf{Trust}: security-aware QEC audit and post-quantum cryptography\\ {\footnotesize RQ5 --- \Cref{ch:qec,ch:pqc}}};
  \node[layer,fill=cbPurple!15,below=of L5] (L4) {\textbf{Synthesis}: learned circuit synthesis toward fault-tolerant budgets; honest cost accounting\\ {\footnotesize RQ4 --- \Cref{ch:synthesis}}};
  \node[layer,fill=cbGreen!15,below=of L4] (L3) {\textbf{Algorithms}: hardware-aware optimization on noisy devices\\ {\footnotesize RQ3 --- \Cref{ch:nisq}}};
  \node[layer,fill=cbOrange!18,below=of L3] (L2) {\textbf{Data}: vectorized encodings for quantum learning\\ {\footnotesize RQ2 --- \Cref{ch:qml}}};
  \node[layer,fill=cbBlue!15,below=of L2] (L1) {\textbf{Scale}: GPU-accelerated circuit simulation on HPC\\ {\footnotesize RQ1 --- \Cref{ch:qgear}}};
  \node[draw,thick,circle,fill=white,minimum size=1.4cm,font=\small\bfseries,right=1.3cm of L3] (qpu) {QPU};
  \draw[arr] (L1.east) -| (qpu.south);
  \draw[arr,<->] (qpu.west) -- (L3.east);
  \draw[arr] (L2.east) -- ++(0.55,0) |- (qpu.220);
  \draw[arr] (qpu.north) |- (L5.east);
  \draw[arr,cbBlue] (L1.west) -- ++(-0.5,0) |- node[pos=0.25,rotate=90,anchor=south,font=\footnotesize] {simulation, containers} (L4.west);
  \draw[arr,cbGray,dashed] (L5.west) -- ++(-1.0,0) |- node[pos=0.25,rotate=90,anchor=north,font=\footnotesize] {seeds, keys} (L4.west);
\end{tikzpicture}
\caption{The HPC-surrounded quantum stack studied in this dissertation. Each classical layer answers one research question and is treated in one or two chapters; artifacts flow between layers as indicated. Schematic.}
\label{fig:intro:stack}
\end{figure}

\section{Contributions}\label{sec:intro:contributions}

The contributions below are grouped by research question.

\paragraph{Scale (RQ1).}
\qgear{} transforms \qiskit{} circuits gate-by-gate into \cudaq{} kernels through a compact three-tensor encoding and packages the simulator as containers for Slurm-managed GPU systems. On the Perlmutter supercomputer it simulates circuits of up to 42 qubits on 1,024 NVIDIA A100 GPUs. It achieves roughly two orders of magnitude speed-up over a 128-core CPU node and roughly an order of magnitude over existing GPU simulators. Users need not rewrite their circuits~\cite{guo2025qgear}.

\paragraph{Data (RQ2).}
Four learning primitives build on vectorized address--data encodings:
\begin{itemize}[nosep,leftmargin=1.5em]
\item \qpie{}, a non-sequential hybrid network with parallel information exchange between classical and quantum branches~\cite{guo2025qpie};
\item \vqt{}, a quantum transformer that computes attention by a vectorized quantum dot product with a learnable nonlinear encoder~\cite{guo2025vqt};
\item \shardq{}, a hardware-aware circuit cutting-and-knitting compiler that reduces encoding error on IBM Heron devices by roughly 15\%~\cite{guo2025shardq}; and
\item \nnqa{}, which compiles trained polynomial networks into exact quantum arithmetic with over 99.5\% accuracy for polynomials up to degree 35 on superconducting and trapped-ion hardware~\cite{guo2026nnqa}.
\end{itemize}

\paragraph{Noise (RQ3).}
\deal{} (direct entanglement ansatz learning) replaces black-box QAOA parameter optimization with a direct mapping from QUBO coefficients to circuit angles. It adds a dynamic physical-qubit mapping driven by device calibration and an adaptive digital zero-noise extrapolation. It improves the ground-state success rate of 20-qubit, 10-layer circuits on IBM Heron processors by up to 14\%~\cite{guo2025deal}. \qawa{} (quantum approximate walk algorithm) introduces a classical-data-traceable oracle. Its depth grows linearly in the number of qubits, and its outputs are interpretable through polynomial-time classical post-processing. \qawa{} demonstrates portfolio optimization with a measurable advantage above the noise floor on a 156-qubit device~\cite{guo2025qawa}.

\paragraph{Synthesis (RQ4).}
\rubriq{} formulates fault-tolerant circuit synthesis as rubric-guided reinforcement learning of a code-generating language model, with GPU-accelerated simulation on Perlmutter inside the reward loop. It achieves a 3.31$\times$ mean T-gate compression and keeps hardware-constraint violations below 1\%~\cite{guo2026rubriq}. A companion study quantifies the measurement and reload costs of direct quantum simulation of nonlinear wave equations with a unified cost-and-error model. It shows that the per-step quantum cost overtakes the classical cost as the grid grows. The coherent nonlinear update is therefore the target that any end-to-end advantage must meet~\cite{guo2026nonlinearwaves}.

\paragraph{Trust (RQ5).}
A bounded cloud fault model and the notion of accepted logical disturbance make the exposure of a fixed public encoder to a co-located fault-injection adversary auditable. A polynomial-cost seeded Clifford ensemble reduces the mean accepted logical disturbance of eight-qubit encoders from 0.150 under a seed-aware adversary to 0.020 under a seed-blind one. The 86.7\% reduction is attributable to rejection. The fixed \code{5}{1}{3} code corrects every tested weight-one fault~\cite{guo2026qecaudit}. The Light Rider QPC SDK is a post-quantum cryptographic module designed to the FIPS 140-3 requirements, with a module-search capability. It provides the validated randomness and key-establishment services that the audit and the rest of the stack consume~\cite{guo2026lightriderqpc}.


\section{Methodology}\label{sec:intro:method}

Three methodological commitments run through the dissertation. The first is to test every algorithmic claim on real quantum hardware and in simulation. The devices used include IBM's Heron-generation superconducting processors (\texttt{ibm\_torino}, \texttt{ibm\_marrakesh}, \texttt{ibm\_kingston}, \texttt{ibm\_pittsburgh}, \texttt{ibm\_boston}), the IBM Nighthawk processor \texttt{ibm\_miami}, and IonQ's Aria-1 and Forte-1 trapped-ion systems. The ideal-simulator baseline always accompanies hardware results, to separate the effect of noise from that of the method.

The second is to treat the classical HPC resources as first-class experimental apparatus. Every chapter reports the GPUs, the number of nodes, and the containers and software versions used. \Cref{app:qgear,app:repro} document the environments and artifacts in enough detail to reproduce them.

The third is honest accounting. When a method delivers no advantage, as in the nonlinear-wave study of \cref{ch:synthesis}, the negative result receives the same care as the positive ones. To know where the crossover lies is itself a research contribution.

\section{Scope and Limitations}\label{sec:intro:scope}

The dissertation is about the software and systems layers of quantum computing, not device physics, control electronics or the design of new error-correcting codes. Its hardware experiments are limited to the qubit counts (up to 36 in the largest trapped-ion run) and depths (a few hundred two-qubit gates) that current cloud-accessible devices support with usable fidelity. The memory of the largest GPU allocation available (1,024 A100 GPUs) limits its simulation experiments and caps exact state-vector simulation at roughly 42--44 qubits.

The security analysis of \cref{ch:qec} uses a bounded fault model with weight-one Pauli faults and small codes. Its extension to correlated faults and to codes at the scale of practical fault tolerance is future work. The cryptographic module of \cref{ch:pqc} is designed to the FIPS 140-3 requirements, but its formal validation status is outside the dissertation's scope. Each chapter discusses its own limitations further, and \cref{ch:conclusion} consolidates them.

\section{Organization of the Dissertation}\label{sec:intro:organization}

\Cref{fig:intro:roadmap} maps each chapter to its research question and its work. \Cref{ch:background} provides the shared background and closes by positioning the dissertation against the literature for each research question. \Cref{ch:qgear,ch:qml,ch:nisq,ch:synthesis} answer RQ1 to RQ4, and \cref{ch:qec,ch:pqc} together answer RQ5. \Cref{ch:conclusion} states what has been established and outlines the research program that follows. Two appendices document the \qgear{} configuration and the research artifacts and reproducibility information.

% alt: A top-to-bottom roadmap: a wide box for the background chapter, a row of five colored chapter boxes below it labelled with their chapter number and research question (the two trust chapters share one box), a wide box for the conclusion and a dashed box for the appendices; arrows run from the background box to each chapter box and from each chapter box to the conclusion.
\begin{figure}[htbp]
\centering
\begin{tikzpicture}[node distance=0.6cm and 0.25cm,
  wide/.style={draw,rounded corners=4pt,fill=white,text width=13.4cm,align=center,font=\small,inner ysep=5pt},
  chap/.style={draw,rounded corners=4pt,text width=2.3cm,align=center,font=\footnotesize\linespread{1}\selectfont,inner ysep=5pt,minimum height=1.4cm},
  arr/.style={-{Stealth[length=2mm]},thick}]
  \node[wide] (bg) {\textbf{\Cref{ch:background}}: Background};
  \node[chap,fill=cbGreen!15,below=of bg] (c3) {\textbf{\Cref{ch:nisq}}\\ RQ3 Noise};
  \node[chap,fill=cbOrange!18,left=of c3] (c2) {\textbf{\Cref{ch:qml}}\\ RQ2 Data};
  \node[chap,fill=cbBlue!15,left=of c2] (c1) {\textbf{\Cref{ch:qgear}}\\ RQ1 Scale};
  \node[chap,fill=cbPurple!15,right=of c3] (c4) {\textbf{\Cref{ch:synthesis}}\\ RQ4 Synthesis};
  \node[chap,fill=cbGray!12,right=of c4] (c5) {\textbf{\Cref{ch:qec,ch:pqc}}\\ RQ5 Trust};
  \node[wide,below=of c3] (concl) {\textbf{\Cref{ch:conclusion}}: Conclusion};
  \node[wide,dashed,below=0.3cm of concl] (app) {Appendices A--B};
  \draw[arr] (c1.north |- bg.south) -- (c1.north);
  \draw[arr] (c2.north |- bg.south) -- (c2.north);
  \draw[arr] (c3.north |- bg.south) -- (c3.north);
  \draw[arr] (c4.north |- bg.south) -- (c4.north);
  \draw[arr] (c5.north |- bg.south) -- (c5.north);
  \draw[arr] (c1.south) -- (c1.south |- concl.north);
  \draw[arr] (c2.south) -- (c2.south |- concl.north);
  \draw[arr] (c3.south) -- (c3.south |- concl.north);
  \draw[arr] (c4.south) -- (c4.south |- concl.north);
  \draw[arr] (c5.south) -- (c5.south |- concl.north);
\end{tikzpicture}
\caption{Roadmap of the dissertation. The background chapter feeds the five chapters that answer RQ1--RQ5, and the conclusion revisits them. Colors follow \cref{fig:intro:stack}.}
\label{fig:intro:roadmap}
\end{figure}
